\begin{lemma}
\label{lemma-argument-proves}
Let $S$ be a scheme. Let $X$ be a quasi-compact and quasi-separated
algebraic space over $S$. Suppose that for every \'etale morphism
$j : V \to W$ with $W \subset X$ quasi-compact open and $V$ affine
the right derived functor
$$
\Phi : D(\QCoh(\mathcal{O}_U)) \to D(\QCoh(\mathcal{O}_W))
$$
of the left exact functor
$j_* : \QCoh(\mathcal{O}_V) \to \QCoh(\mathcal{O}_W)$
fits into a commutative diagram
$$
\xymatrix{
D(\QCoh(\mathcal{O}_V)) \ar[d]_\Phi \ar[r]_{i_V} &
D_\QCoh(\mathcal{O}_V) \ar[d]^{Rj_*} \\
D(\QCoh(\mathcal{O}_W)) \ar[r]^{i_W} &
D_\QCoh(\mathcal{O}_W)
}
$$
Then the functor (\ref{equation-compare})
$$
D(\QCoh(\mathcal{O}_X))
\longrightarrow
D_\QCoh(\mathcal{O}_X)
$$
is an equivalence with quasi-inverse given by $RQ_X$.
\end{lemma}

\begin{proof}
We first use the induction principle to prove $i_X$ is fully faithful.
More precisely, we will use Lemma \ref{lemma-induction-principle-enlarge}.
Let $(U \subset W, V \to W)$ be an elementary distinguished square
with $V$ affine and $U, W$ quasi-compact open in $X$. Assume that
$i_U$ is fully faithful. We have to show that $i_W$ is fully faithful.
We may replace $X$ by $W$, i.e., we may assume $W = X$ (we do this just
to simplify the notation -- observe that the condition in the
statement of the lemma is preserved under this operation).

\medskip\noindent
Suppose that $A, B$ are objects of $D(\QCoh(\mathcal{O}_X))$.
We want to show that
$$
\Hom_{D(\QCoh(\mathcal{O}_X))}(A, B)
\longrightarrow
\Hom_{D(\mathcal{O}_X)}(i_X(A), i_X(B))
$$
is bijective. Let $T = |X| \setminus |U|$.

\medskip\noindent
Assume first $i_X(B)$ is supported on $T$. In this case the map
$$
i_X(B) \to Rj_{V, *}(i_X(B)|_V) = Rj_{V, *}(i_V(B|_V))
$$
is a quasi-isomorphism
(Lemma \ref{lemma-pushforward-with-support-in-open}).
By assumption we have an isomorphism
$i_X(\Phi(B|_V)) \to Rj_{V, *}(i_V(B|_V))$ in $D(\mathcal{O}_X)$.
Moreover, $\Phi$ and ${-}|_V$ are adjoint functors on the derived categories of
quasi-coherent modules (by
Derived Categories, Lemma \ref{derived-lemma-derived-adjoint-functors}).
The adjunction map $B \to \Phi(B|_V)$ becomes an isomorphism
after applying $i_X$, whence is an isomorphism in
$D(\QCoh(\mathcal{O}_X))$.
Hence
\begin{align*}
\Mor_{D(\QCoh(\mathcal{O}_X))}(A, B)
& =
\Mor_{D(\QCoh(\mathcal{O}_X))}(A, \Phi(B|_V)) \\
& =
\Mor_{D(\QCoh(\mathcal{O}_V))}(A|_V, B|_V) \\
& =
\Mor_{D(\mathcal{O}_V)}(i_V(A|_V), i_V(B|_V)) \\
& =
\Mor_{D(\mathcal{O}_X)}(i_X(A), Rj_{V, *}(i_V(B|_V))) \\
& =
\Mor_{D(\mathcal{O}_X)}(i_X(A), i_X(B))
\end{align*}
as desired. Here we have used that $i_V$ is fully faithful
(Lemma \ref{lemma-affine-coherator}).

\medskip\noindent
In general, choose any complex $\mathcal{B}^\bullet$ of quasi-coherent
$\mathcal{O}_X$-modules representing $B$. Next, choose any quasi-isomorphism
$s : \mathcal{B}^\bullet|_U \to \mathcal{C}^\bullet$ of complexes of
quasi-coherent modules on $U$. As $j_U : U \to X$ is
quasi-compact and quasi-separated the functor $j_{U, *}$ transforms
quasi-coherent modules into quasi-coherent modules
(Morphisms of Spaces, Lemma \ref{spaces-morphisms-lemma-pushforward}).
Thus there is a canonical map
$\mathcal{B}^\bullet \to j_{U, *}(\mathcal{B}^\bullet|_U) \to
j_{U, *}\mathcal{C}^\bullet$
of complexes of quasi-coherent modules on $X$.
Set $B'' = j_{U, *}\mathcal{C}^\bullet$ in $D(\QCoh(\mathcal{O}_X))$
and choose a distinguished triangle
$$
B \to B'' \to B' \to B[1]
$$
in $D(\QCoh(\mathcal{O}_X))$. Since the first arrow of the triangle
restricts to an isomorphism over $U$ we see that $B'$ is supported on $T$.
Hence in the diagram
$$
\xymatrix{
\Hom_{D(\QCoh(\mathcal{O}_X))}(A, B'[-1]) \ar[r] \ar[d] &
\Hom_{D(\mathcal{O}_X)}(i_X(A), i_X(B')[-1]) \ar[d] \\
\Hom_{D(\QCoh(\mathcal{O}_X))}(A, B) \ar[r] \ar[d] &
\Hom_{D(\mathcal{O}_X)}(i_X(A), i_X(B)) \ar[d] \\
\Hom_{D(\QCoh(\mathcal{O}_X))}(A, B'') \ar[r] \ar[d] &
\Hom_{D(\mathcal{O}_X)}(i_X(A), i_X(B'')) \ar[d] \\
\Hom_{D(\QCoh(\mathcal{O}_X))}(A, B') \ar[r] &
\Hom_{D(\mathcal{O}_X)}(i_X(A), i_X(B'))
}
$$
we have exact columns and the top and bottom horizontal arrows are
bijective. Finally, choose a complex $\mathcal{A}^\bullet$
of quasi-coherent modules representing $A$.

\medskip\noindent
Let $\alpha : i_X(A) \to i_X(B)$ be a morphism between
in $D(\mathcal{O}_X)$. The restriction $\alpha|_U$ comes from a
morphism in $D(\QCoh(\mathcal{O}_U))$ as $i_U$ is fully faithful.
Hence there exists a choice of
$s : \mathcal{B}^\bullet|_U \to \mathcal{C}^\bullet$ as above
such that $\alpha|_U$ is represented by an actual map of complexes
$\mathcal{A}^\bullet|_U \to \mathcal{C}^\bullet$.
This corresponds to a map of complexes
$\mathcal{A} \to j_{U, *}\mathcal{C}^\bullet$.
In other words, the image of $\alpha$ in
$\Hom_{D(\mathcal{O}_X)}(i_X(A), i_X(B''))$ comes from
an element of $\Hom_{D(\QCoh(\mathcal{O}_X))}(A, B'')$.
A diagram chase then shows that $\alpha$ comes from a morphism
$A \to B$ in $D(\QCoh(\mathcal{O}_X))$. Finally, suppose
that $a : A \to B$ is a morphism of $D(\QCoh(\mathcal{O}_X))$
which becomes zero in $D(\mathcal{O}_X)$. After choosing $\mathcal{B}^\bullet$
suitably, we may assume $a$ is represented by a morphism of complexes
$a^\bullet : \mathcal{A}^\bullet \to \mathcal{B}^\bullet$.
Since $i_U$ is fully faithul the restriction $a^\bullet|_U$ is zero
in $D(\QCoh(\mathcal{O}_U))$. Thus we can choose $s$
such that
$s \circ a^\bullet|_U : \mathcal{A}^\bullet|_U \to \mathcal{C}^\bullet$
is homotopic to zero. Applying the functor $j_{U, *}$ we conclude that
$\mathcal{A}^\bullet \to j_{U, *}\mathcal{C}^\bullet$ is homotopic
to zero. Thus $a$ maps to zero in
$\Hom_{D(\QCoh(\mathcal{O}_X))}(A, B'')$.
Thus we may assume that $a$ is the image of an element
of $b \in \Hom_{D(\QCoh(\mathcal{O}_X))}(A, B'[-1])$.
The image of $b$ in $\Hom_{D(\mathcal{O}_X)}(i_X(A), i_X(B')[-1])$
comes from a $\gamma \in \Hom_{D(\mathcal{O}_X)}(A, B''[-1])$
(as $a$ maps to zero in the group on the right). Since we've
seen above the horizontal arrows are surjective, we see
that $\gamma$ comes from a $c$ in
$\Hom_{D(\QCoh(\mathcal{O}_X))}(A, B''[-1])$
which implies $a = 0$ as desired.

\medskip\noindent
At this point we know that $i_X$ is fully faithful for our original $X$.
Since $RQ_X$ is its right adjoint, we see that
$RQ_X \circ i_X = \text{id}$ (Categories, Lemma
\ref{categories-lemma-adjoint-fully-faithful}).
To finish the proof we show that for any
$E$ in $D_\QCoh(\mathcal{O}_X)$ the map
$i_X(RQ_X(E)) \to E$ is an isomorphism. Choose a distinguished triangle
$$
i_X(RQ_X(E)) \to E \to E' \to i_X(RQ_X(E))[1]
$$
in $D_\QCoh(\mathcal{O}_X)$. A formal argument using the
above shows that $i_X(RQ_X(E')) = 0$. Thus it suffices to prove that
for $E \in D_\QCoh(\mathcal{O}_X)$ the condition
$i_X(RQ_X(E)) = 0$ implies that $E = 0$. Consider an \'etale morphism
$j : V \to X$ with $V$ affine. By
Lemmas \ref{lemma-affine-coherator} and
\ref{lemma-flat-pushforward-coherator}
and our assumption we have
$$
Rj_*(E|_V) = Rj_*(i_V(RQ_V(E|_V))) = i_X(\Phi(RQ_V(E|_V))) =
i_X(RQ_X(Rj_*(E|_V)))
$$
Choose a distinguished triangle
$$
E \to Rj_*(E|_V) \to E' \to E[1]
$$
Apply $RQ_X$ to get a distinguished triangle
$$
0 \to RQ_X(Rj_*(E|_V)) \to RQ_X(E') \to 0[1]
$$
in other words the map in the middle is an isomorphism.
Combined with the string of equalities above we find
that our first distinguished triangle becomes a distinguished triangle
$$
E \to i_X(RQ_X(E')) \to E' \to E[1]
$$
where the middle morphism is the adjunction map. However, the composition
$E \to E'$ is zero, hence $E \to i_X(RQ_X(E'))$ is zero by adjunction!
Since this morphism is isomorphic to the morphism
$E \to Rj_*(E|_V)$ adjoint to $\text{id} : E|_V \to E|_V$ we
conclude that $E|_V$ is zero. Since this holds for all
affine $V$ \'etale over $X$ we conclude $E$ is zero as desired.
\end{proof}